\section{Conclusion}
\label{sec:conclusion}

This paper studied production--logistics collaborative scheduling with ten
modelled constraints and three reported objectives, using a joint-chain
formulation in which all production and logistics decisions of an episode form
a single decision chain trained by group-relative policy optimization. The
policy's parts are standard, a Transformer encoder with candidate-scoring
heads; what the paper contributes is the chain that joins the two decision
layers, the encoding that puts the aisle network into the state, and the
measurement of what the ten constraints do when they are all present.

The method beats a rule baseline by 26.5 in makespan on the fully constrained
problem and dominates every out-of-sample point of an NSGA-II front, while
losing narrowly to the same metaheuristic on the unconstrained problem. That
combination locates the method's advantage in its representation of the
constrained state rather than in its search behaviour. On the objective side,
the aggregate energy consumption commonly reported in this literature is close
to rank-equivalent with the makespan, because the terms that dominate it are
linear in the horizon by construction; under that account the reported
three-objective problem has an effective dimension of two, which is consistent
with the observed collapse of the preference front to two points. A net account
that removes the horizon-linear terms restores independent ordering.

\paragraph{Limitations} The evaluation is a first pass and the scope is narrow.
Every number comes from a single instance of the benchmark, and MK01 is the
instance least favourable to the batching mechanism, so the batching results
here are not a verdict on batching. Except for the advantage ablation, each
configuration was trained once. Against that, the training-seed standard
deviation is measured rather than assumed: it is 5.90 makespan units on the same
configuration, which is the order of most of the differences we report, and it
is why the ablation table is presented without a ranking.

The comparison uses a rule, a single-objective version of our own method, and
NSGA-II. It does not use a published deep reinforcement learning method, which
is the baseline a reviewer would most reasonably ask for; adapting one to this
simulation is substantial work, and Section~\ref{sec:related:incomparable}
explains the boundary conditions under which any such number could be compared.
A state-of-the-art metaheuristic for the flexible job shop is absent for the
same reason. On the constraint side, two of the ten do not activate under the
reported parameters: the charging constraint never fires across the benchmark,
and maintenance does not trigger on MK01, so the budget mechanism penalizes four
constraints in practice rather than the six it is configured with. The
consequence is stated plainly: this paper reports no evidence on whether the
charging mechanism works, and establishes that only in a configuration where
vehicles actually drain. Most of the constraint parameters have no published
source and are values we chose; the sensitivity analysis in
Section~\ref{sec:results:sensitivity} reports how far the results move when they
change, but does not make them cited. One parameterization of the reported
energy is conservative: the machine standby power is taken from the loaded-idle
column of the cited source rather than its separately tabulated standby power,
which overstates standby energy by 12.9\% of the total without changing any
ordering. Finally, the paper proves no bound and no convergence guarantee.

\paragraph{Future work} Three measurements would change the conclusions rather
than extend them. The ablation table reports a grouping, not an ordering,
because one training seed per arm cannot separate differences of the size it
reports. An ordering would need about ten seeds per arm, some hundred and fifty
runs at the cost measured in Section~\ref{sec:method:cost}, with the number set
by the smallest difference the table is meant to resolve: at a seed-to-seed
spread of eight makespan units, an effect of eight units needs eight seeds and
an effect of five needs twenty. The comparison against the baselines is much
cheaper to settle, because its effect is three times the seed-to-seed spread;
three seeds per arm would certify it. Running the mechanisms on further
instances of the benchmark would separate the effects that belong to the
mechanism from those that belong to MK01, which matters most for batching and
for maintenance. And training with the net energy account as the reward, rather
than measuring it after the fact, would test whether the third dimension
available at the measurement level is usable at the optimization level. The
comparison against a published deep reinforcement learning method remains the
piece of work this paper most needs done next.
